% Exact follow-on module. Input after the original Niemeier and Leech proofs.
% Original Euclidean coordinates, residue action, scales and glue are retained.
\section{The residue cubic action and the common integral Leech correspondence}
\label{esrf:sec:fixed-neighbor}

Retain the exact lattices $R=A^{\perp4}\perp D\subset N$ and the
binary and ternary codes from the index-72 construction.
The four $A$ coordinates remain the ordered six-element residue
fibres. To avoid confusing a vector with a representation, write
$\varrho$ for the Weyl vector
\[
 \varrho_A=(5,3,1,-1,-3,-5)/2,\qquad
 \varrho_D=(3,2,1,0),\qquad
 \varrho=\varrho_A^{\oplus4}\perp\varrho_D.
\]
Thus $\varrho\in N$ and $\varrho^2=84$.
The original residue representation is still denoted by $\rho$.
The specific element studied here is
\[
 r=\rho(361)=\rho(289)^2,\qquad
 r|_{A_i}=C^2,\quad Ce_j=e_{j+1\bmod6},\qquad r|_D=I.
 \tag{RF1}
\]
In particular $r^3=I$. This is the actual residue coordinate
permutation, with no conjugation or replacement.

For each $\sigma\in\Omega=O(N)\varrho$, choose an original root
$\alpha_\sigma$ with $(\sigma,\alpha_\sigma)=1$ and set
\[
 \begin{split}
 \ell_\sigma(x)&=(\sigma,x)\bmod6,\qquad
 K_\sigma=\ker(\ell_\sigma:N\to\mathbb Z/6),\\
 \Lambda_\sigma&=K_\sigma+
                      \mathbb Z(\sigma/6-\alpha_\sigma)\\
 &=\{x+k\sigma/6:x\in N,\ k\in\mathbb Z,\
                                    (\sigma,x)+k\equiv0\pmod6\}.
 \end{split}
 \tag{RF2}
\]
The preceding neighbor proof shows that the definition is independent
of $\alpha_\sigma$, every $\Lambda_\sigma$ is the specified rootless
even unimodular lattice, and
$N\cap\Lambda_\sigma=K_\sigma$ with both indices six.
The family considered in this section is precisely
$\mathscr F=\{\Lambda_\sigma:\sigma\in\Omega\}$.

\subsection{A necessary and sufficient equality criterion}

\begin{theorem}[Equality of the marked Weyl neighbors]
\label{esrf:thm:equality}
For $\sigma,\tau\in\Omega$, the following three statements are equivalent:
\[
 \Lambda_\sigma=\Lambda_\tau,\qquad
 K_\sigma=K_\tau,\qquad
 \tau-\epsilon\sigma\in6N
                 \text{ for some }\epsilon\in\{1,-1\}.
 \tag{RF3}
\]
\end{theorem}

\begin{proof}
Equality of the two neighbors gives equality of the kernels by
intersection with $N$.
Each $\ell_\sigma$ is onto because
$(\sigma,\alpha_\sigma)=1$.
If the kernels are equal, their induced isomorphisms
$N/K_\sigma\to\mathbb Z/6$ differ by an automorphism of
$\mathbb Z/6$. Its two units are $1,-1$.
Therefore $(\tau-\epsilon\sigma,N)\subset6\mathbb Z$.
Unimodularity $N^\vee=N$ gives
$(\tau-\epsilon\sigma)/6\in N$.
Conversely this congruence makes the two kernels equal.

It remains to show that this kernel criterion actually gives
equality of the chosen unimodular extensions.
The exact formula (RF2), after replacing $k$ by $-k$, gives
$\Lambda_{-\sigma}=\Lambda_\sigma$.
We may consequently suppose $\tau=\sigma+6t$ with $t\in N$.
Both Weyl vectors have squared norm 84, so
\[
 0=\tau^2-\sigma^2=12(\sigma,t)+36t^2,\qquad
                       (\sigma,t)=-3t^2\in6\mathbb Z.
 \tag{RF4}
\]
The last inclusion uses the actual evenness of $N$.
For $x+k\tau/6$ in (RF2), put $x'=x+kt\in N$.
Then $x+k\tau/6=x'+k\sigma/6$, and
\[
 (\sigma,x')+k
 \equiv(\sigma,x)+k
 \equiv(\tau,x)+k\equiv0\pmod6.
\]
This proves $\Lambda_\tau\subset\Lambda_\sigma$.
Interchanging $\sigma,\tau$ proves the reverse inclusion.
Thus no additional unverified choice of an extension is being
identified merely from equality of its common kernel.
\end{proof}

It follows in particular that the map
\[
 \Omega/{\sim}\ \longrightarrow\ \mathscr F,\qquad
 [\sigma]\longmapsto\Lambda_\sigma,\qquad
 \sigma\sim\tau\ \Longleftrightarrow\
                          \tau\equiv\pm\sigma\pmod{6N},
 \tag{RF5}
\]
is a bijection. It is $O(N)$-equivariant: the original formula
$T\Lambda_\sigma=\Lambda_{T\sigma}$ follows by substitution in
(RF2), and $T$ preserves the congruence relation since $T(N)=N$.

\subsection{The exact tetracode obstruction to a fixed neighbor}

\begin{theorem}[No fixed member of this specified family]
\label{esrf:thm:no-fixed-neighbor}
The actual residue element $r=\rho(361)$ fixes no lattice of
$\mathscr F$. Every orbit of its action on $\mathscr F$ has
exactly three members.
\end{theorem}

\begin{proof}
Suppose first that $r\sigma\equiv-\sigma\pmod{6N}$.
Applying $r$ successively, using $r(N)=N$ and $r^3=I$, gives
$\sigma\equiv-\sigma\pmod{6N}$.
Then $\sigma/3\in N$, contradicting its pairing
$1/3$ with the root $\alpha_\sigma\in N$.
Thus (RF3) permits fixedness only if
$t=(r\sigma-\sigma)/6\in N$.

Every isometry of $N$ preserves the original root lattice, permutes
its four $A_5$ components, and acts in each of them by a signed
coordinate permutation, as proved in the full automorphism theorem.
Negation permutes the six entries of $\varrho_A$.
Consequently each $A$ block $s=(s_0,\ldots,s_5)$ of every
$\sigma\in\Omega$ is a permutation of the six distinct numbers
$5/2,3/2,1/2,-1/2,-3/2,-5/2$.
In this block put $\Delta_j=s_{j-2\bmod6}-s_j$.
If the corresponding block $\Delta/6$ of $t$ belongs to $A^\vee$,
its coordinate differences are integral, so
\[
 \Delta_j\equiv c\pmod6\qquad\text{for every }j
 \tag{RF6}
\]
for one integer $c$ modulo six.
On either three-cycle of the index permutation, the three
$\Delta_j$ sum to zero. Hence $3c\equiv0\pmod6$ and
$c\in\{0,2,4\}$.
Each $\Delta_j$ is an integer between $-5$ and $5$.
If $c=0$, (RF6) forces all $\Delta_j=0$, contrary to the
distinctness of the entries on each three-cycle.
Therefore $c=2$ or $4$.

In the original discriminant marking, a vector of the $A$ class
$t_i\lambda+A$ has common fractional part $-t_i/6$.
The corresponding class of $\Delta/6$ is therefore
$t_i=-c\pmod6$, also in $\{2,4\}$.
The difference on the $D$ block is exactly zero because $r|_D=I$.
Thus membership of the full vector $(r\sigma-\sigma)/6$ in $N$
would require a glue word
\[
 (t_1,t_2,t_3,t_4;00)\in H,\qquad t_i\in\{2,4\}
                                                \text{ for every }i.
 \tag{RF7}
\]
In $H=\{(3u+2z;y)\}$ this forces $u=0$ and all four $z_i$
nonzero. But the exact ternary code is
\[
 C_3=\{(a+b,a+2b,a,b):a,b\in\mathbb F_3\}.
 \tag{RF8}
\]
If exactly one parameter is zero, exactly one coordinate is zero.
If both are nonzero, precisely one of $a+b,a+2b$ is zero.
Thus every nonzero word has weight three, and the zero word
has weight zero. No word in (RF7) belongs to $H$.
This contradicts $t\in N$.

Neither sign in (RF3) can occur. Since $r^3=I$, every orbit
therefore has cardinality three.
\end{proof}

The conclusion concerns exactly $\mathscr F$ and the action (RF1).
The integral correspondence for these three-member orbits is
constructed next; no assertion about other embeddings or
constructions of a Leech lattice follows from this obstruction.

\subsection{The canonical orbit and its common integral lattice}

Put $\sigma_j=r^j\varrho$, $\Lambda_j=\Lambda_{\sigma_j}$ and
$K_j=K_{\sigma_j}$, with $j$ read modulo three.
Keep the same original root
$\alpha=e_0-e_1$ in the first $A$ block.
Directly $(\sigma_j,\alpha)=1$ for all three $j$.
Define
\[
 v_j=\sigma_j/6-\alpha,\quad v=v_0,\qquad
 q_1=(\sigma_1-\sigma_0)/6,\quad
 q_2=(\sigma_2-\sigma_0)/6.
 \tag{RF9}
\]
The complete original coordinates are
\[
 \begin{split}
 q_1&=\bigl((-2,-2,1,1,1,1)/3\bigr)^{\oplus4}\perp0,\\
 q_2&=\bigl((-1,-1,-1,-1,2,2)/3\bigr)^{\oplus4}\perp0.
 \end{split}
 \tag{RF10}
\]
Their classes in $R^\vee/R$ are respectively
$(4,4,4,4;00)$ and $(2,2,2,2;00)$.
Neither is in $H$, while $3q_1,3q_2\in R$.
In particular each $q_i+N$ has exact order three.
Direct pairings give
\[
 q_1^2=q_2^2=16/3,\quad(q_1,q_2)=8/3,\quad
                    (\varrho,q_1)=(\varrho,q_2)=-16,
 \qquad q_1+q_2\in R.
 \tag{RF11}
\]

For $x\in N$, write its original glue class as $(3u+2z;y)$,
with $z=(a+b,a+2b,a,b)\in C_3$, and define the homomorphism
\[
 \beta:N\longrightarrow\mathbb F_3,\qquad\beta(x)=b=z_4.
 \tag{RF12}
\]
This map is defined through the original quotient $N/R$.
This definition uses the retained residue marking without changing it.
The sum of the four $z$ coordinates is also $b$ modulo three.
Write $\ell_j(x)=(\sigma_j,x)\bmod6$.
Then the exact three character formulas are
\[
 \ell_1(x)=\ell_0(x)+4\beta(x),\qquad
 \ell_2(x)=\ell_0(x)+2\beta(x)\quad\text{in }\mathbb Z/6.
 \tag{RF13}
\]
The maps $b\mapsto2b,4b$ from $\mathbb F_3$ to $\mathbb Z/6$
are well-defined homomorphisms.
To prove (RF13), the discriminant bilinear form gives
\[
 (q_1,x)\equiv\tfrac56\,4\sum t_i
       \equiv\tfrac13\sum t_i
       \equiv\tfrac23\beta(x)\pmod{\mathbb Z},
 \qquad t_i=3u_i+2z_i\pmod6.
\]
For $q_2$ the same calculation gives
$(q_2,x)\equiv\beta(x)/3\pmod{\mathbb Z}$.
Multiplication by six and (RF9) prove both formulas.

Using the original dominant representatives
$\omega_t=(1^t,0^{6-t})-\frac t6(1,\ldots,1)$, put
\[
 u_*=(\omega_2,\omega_4,0,\omega_2;0)\in N.
 \tag{RF14}
\]
Its ternary class is exactly the original generator
$z_2=(1,2,0,1)$.
It has $u_*^2=4$, $(\varrho,u_*)=12$ and $\beta(u_*)=1$.
The map
\[
 (\ell_0,\beta):N\longrightarrow\mathbb Z/6\oplus\mathbb F_3
 \tag{RF15}
\]
is onto: $\alpha$ maps to $(1,0)$ and $u_*$ to $(0,1)$.
Consequently
\[
 \begin{gathered}
 I=\ker(\ell_0,\beta)=K_0\cap K_1=K_1\cap K_2
     =K_2\cap K_0=\bigcap_jK_j,\\
 [N:I]=18,\qquad[K_j:I]=3,\qquad\det I=324.
 \end{gathered}
 \tag{RF16}
\]
The kernel equalities follow from (RF13), because either
$2b=0$ or $4b=0$ in $\mathbb Z/6$ forces $b=0$ in
$\mathbb F_3$.
The action $r$ permutes the $K_j$, so it preserves $I$.
Also $3q_i\in I$: these are in $R$, hence have $\beta=0$,
and $(\varrho,3q_i)=-48$ is divisible by six.

Define the full-rank lattice
\[
 J=I+\mathbb Z(3v).
 \tag{RF17}
\]
It is contained in $\Lambda_0$.
The element $6v=\varrho-6\alpha$ lies in $I$: its
$\ell_0$ value is zero, and its glue has zero ternary part.
The element $3v$ is not in $N$, because
$(3v,\alpha)=-11/2$ is not an integer.
Thus
\[
 N\cap J=I,\qquad[J:I]=2,\qquad
                    \det J=324/4=81,\qquad[\Lambda_0:J]=9.
 \tag{RF18}
\]
This lattice is positive definite, even, and rootless, because
it is a sublattice of $\Lambda_0$.

All exact root corrections in the cyclic action are retained.
Put $\kappa_j=\alpha-r^j\alpha\in R$.
The vectors $r\alpha=e_2-e_3$ and $r^2\alpha=e_4-e_5$
are original simple roots with $\varrho$-pairing one.
Hence $\kappa_j\in I$.
The precise formulas, rather than literal identification of
different generators, are
\[
 r^jv=v+q_j+\kappa_j\quad(j=1,2),\qquad
 r(3v)=3v+3q_1+3\kappa_1.
 \tag{RF19}
\]
Since $rI=I$, $3q_1\in I$ and $\kappa_1\in I$, this proves
$rJ\subset J$. Applying $r^3=I$ gives equality $rJ=J$.

\begin{theorem}[All three pairwise Leech intersections]
\label{esrf:thm:common-J}
The actual three-member orbit satisfies
\[
 \Lambda_0\cap\Lambda_1
 =\Lambda_1\cap\Lambda_2
 =\Lambda_2\cap\Lambda_0
 =\bigcap_j\Lambda_j=J,\qquad[\Lambda_j:J]=9.
 \tag{RF20}
\]
The element $r$ is an integral isometry of $J$, and induces
isometries $\Lambda_j\to\Lambda_{j+1}$ with this common restriction.
\end{theorem}

\begin{proof}
First $J\subset\Lambda_0\cap\Lambda_1$: $I$ lies in both,
and $3v=3v_1-3q_1$ with $3q_1\in I$.
Suppose $y\in\Lambda_0\cap\Lambda_1$.
Using the same $\alpha$, write
$y=x_0+k_0v_0=x_1+k_1v_1$, $x_j\in K_j$.
The fractional part of $(v_j,\alpha)$ is $1/6$ for both
$j$, so pairing with $\alpha\in N$ gives
$k_0\equiv k_1\pmod6$.
Absorbing multiples of $6v_j\in K_j$, take the same
$k\in\{0,\ldots,5\}$ in both expressions.
Then $x_0-x_1=kq_1\in N$.
The exact order three of $q_1+N$ forces $k=0$ or $3$.
For these values $kq_1\in I\subset K_1$, so
$x_0=x_1+kq_1\in K_0\cap K_1=I$.
Thus $y=x_0+kv\in J$, proving the first equality.
Apply $r$ and $r^2$, using $rJ=J$, to prove the other pairwise
equalities. Their common index is nine by (RF18) and the
isometries induced by $r$.
\end{proof}

For comparison with the original index-72 glue, let
$R_0=R\cap K_0$. Since $\beta|_R=0$, (RF13) gives
$R\cap K_j=R_0$ for every $j$, and $R\cap I=R_0$.
The reduction map $K_0/R_0\cong H$ identifies $I/R_0$
with the kernel of the original ternary coefficient $b$.
Thus
\[
 [I:R_0]=24,\qquad[J:R_0]=48,\qquad
 [K_0:R_0]=72,\qquad[\Lambda_j:R_0]=432.
 \tag{RF21}
\]
All these are inclusions of the displayed original lattices.

\subsection{The exact discriminant form and its cyclic three planes}

\begin{theorem}[The integral correspondence over $\mathbb F_3^4$]
\label{esrf:thm:discriminant-correspondence}
One has
\[
 J^\vee=\Lambda_i+\Lambda_j\quad(i\ne j),\qquad
 J^\vee/J\cong\mathbb F_3^4.
 \tag{RF22}
\]
In the ordered basis
\[
 e_1=u_*+J,\quad e_2=v+J,\quad
 e_3=ru_*+J,\quad e_4=rv+J,
 \tag{RF23}
\]
the discriminant quadratic form is
\[
 q_J(a,b,c,d)=\frac23(2ad+bc+bd)\pmod{2\mathbb Z},
 \tag{RF24}
\]
and three times its discriminant bilinear form has matrix
\[
 B=
 \begin{pmatrix}
 0&0&0&2\\0&0&1&1\\0&1&0&0\\2&1&0&0
 \end{pmatrix}\quad\text{over }\mathbb F_3.
 \tag{RF25}
\]
The three original Leech extensions are exactly the inverse
images in $J^\vee$ of the isotropic planes
\[
 \begin{split}
 P_0&=\{(a,b,0,0):a,b\in\mathbb F_3\},\\
 P_1&=\{(0,0,c,d):c,d\in\mathbb F_3\},\\
 P_2&=\{(a,b,a-b,b):a,b\in\mathbb F_3\}.
 \end{split}
 \tag{RF26}
\]
The induced action of the actual $r$ is the matrix
\[
 \mathcal R=
 \begin{pmatrix}
 0&0&2&1\\
 0&0&0&2\\
 1&0&2&2\\
 0&1&0&2
 \end{pmatrix},\qquad
 \mathcal R^3=I,\quad \mathcal R^TB\mathcal R=B,
 \quad\mathcal R P_j=P_{j+1}.
 \tag{RF27}
\]
All entries in (RF23)--(RF27) use the original glue generator
$z_2$, root $\alpha$, pairing, and residue permutation.
\end{theorem}

\begin{proof}
Since $J\subset\Lambda_i=\Lambda_i^\vee$, each $\Lambda_i$
lies in $J^\vee$.
By (RF20), the two subgroups $\Lambda_i/J$ and $\Lambda_j/J$
have order nine and intersect trivially for $i\ne j$.
Their sum has order 81, equal to
$|J^\vee/J|=\det J=81$, proving the first equality in (RF22).

The homomorphism $\beta:K_0/I\to\mathbb F_3$ is an isomorphism,
with generator $u_*+I$. Hence $\Lambda_0/J$ is generated by
$u_*+J,v+J$.
Both have order dividing three because $3u_*\in I$ and
$3v\in J$.
They are independent: if $a u_*+b v\in J$, reduce first modulo
$N$; the subgroup $J+N$ has the two cosets represented by
$0,3v$, whereas $v+N$ has order six.
For $a,b\in\{0,1,2\}$ this forces $b=0$.
Then membership of $a u_*$ in $I=N\cap J$ forces $a=0$.
Thus $\Lambda_0/J\cong\mathbb F_3^2$ with the first two
generators in (RF23).
Applying $r$ proves the corresponding statement for
$\Lambda_1/J$ and the last two generators.
Their trivial intersection and total order 81 prove that
(RF23) is indeed a basis of the full discriminant group.

Here are the full Euclidean Gram entries before any passage
to a quotient:
\[
 \operatorname{Gram}(u_*,v,ru_*,rv)=
 \begin{pmatrix}
 4&2&-2&-4/3\\
 2&4&-2/3&-2/3\\
 -2&-2/3&4&2\\
 -4/3&-2/3&2&4
 \end{pmatrix}.
 \tag{RF28}
\]
For an explicit verification, each of the three nonzero
$A$ blocks of $u_*$ pairs to $-2/3$ with its shifted block,
giving $(u_*,ru_*)=-2$.
One has
$(u_*,\alpha)=(u_*,r\alpha)=(ru_*,\alpha)=0$,
$(\varrho,u_*)=12$,
$(\sigma_1,u_*)=-8$, and
$(\sigma_2,u_*)=-4$, directly from (RF10) and (RF14).
These give respectively
$(u_*,v)=2$, $(u_*,rv)=-4/3$, and
$(v,ru_*)=-2/3$.
Also $(\varrho,\sigma_1)=-12$,
$(\varrho,r\alpha)=(\sigma_1,\alpha)=1$,
and $(\alpha,r\alpha)=0$. Thus
\[
 (v,rv)=-12/36-1/6-1/6=-2/3.
\]
The four diagonal entries equal four, and
$(ru_*,rv)=(u_*,v)=2$ by the exact isometry $r$.
This proves (RF28).
Reducing its entries modulo $\mathbb Z$ gives (RF25).
The diagonal and integral cross terms in the squared norm
are even. Reducing the remaining cross terms modulo
$2\mathbb Z$ gives exactly (RF24).
Changing any coordinate by three changes (RF24) by an even
integer, so it is well-defined on the asserted field basis.

It remains to calculate the action without identifying a
family permutation with a fixed Leech automorphism.
By definition $r e_1=e_3$, $r e_2=e_4$.
The vector $u_*+ru_*+r^2u_*$ lies in $I$: its $\beta$ value
is $1+1+1=0$, because $r$ acts trivially on $N/R$,
and its $\ell_0$ value is $12-4-8=0$.
Consequently
\[
 r^2u_*\equiv-u_*-ru_*\pmod J.
 \tag{RF29}
\]
By (RF19),
\[
 v+rv+r^2v
       =3v+q_1+q_2+\kappa_1+\kappa_2.
\]
The $\kappa_j$ belong to $I$, and $3v\in J$.
Furthermore $q_1+q_2-u_*+ru_*\in I$.
Indeed $q_1+q_2\in R$, so its $\beta$ value is zero;
subtracting $u_*$ and adding $ru_*$ keeps that value zero.
Its pairing with $\varrho$ is
$-32-12-4=-48$, divisible by six.
Thus
\[
 r^2v\equiv u_*-v-ru_*-rv\pmod J.
 \tag{RF30}
\]
The four images in (RF29)--(RF30) give the four columns of
$\mathcal R$ in (RF27).
The matrix has cube the identity and preserves $B$ either
by multiplication, or by the already proved facts $r^3=I$
and $rJ=J$ with its original isometric pairing.

The first two planes in (RF26) are respectively
$\Lambda_0/J$ and $\Lambda_1/J$ by their proved bases.
Their next image is the span of
$(-1,0,-1,0)$ and $(1,-1,-1,-1)$, from (RF29)--(RF30).
Solving their two linear parameters gives exactly
$(a,b,a-b,b)$, proving the stated $P_2=\Lambda_2/J$.
Each plane is isotropic: the first two vanish in (RF24),
and the third has
$2ab+b(a-b)+b^2=3ab=0$ in $\mathbb F_3$.
They have nine elements and pairwise intersection zero,
also following directly from (RF26).
Finally each $\Lambda_j$ is exactly the inverse image of
its quotient plane because it contains $J$.
This proves the claimed integral correspondence and all
its maps, indices, pairings, and cyclic actions.
\end{proof}

\subsection{The fixed fourth plane and its exact \texorpdfstring{$A_2^{12}$}{A2 to the twelfth power} extension}

The same finite discriminant form has the additional plane
\[
 P_\infty=\{(a,b,-a+b,-b):a,b\in\mathbb F_3\}.
 \tag{RF31}
\]
Put
\[
 p=u_*-ru_*,\qquad q=v+ru_*-rv,\qquad
                       M=J+\mathbb Zp+\mathbb Zq.
 \tag{RF32}
\]
These are literal vectors in the original Euclidean space.

\begin{theorem}[The fixed fourth extension]
\label{esrf:thm:fourth-extension}
The plane $P_\infty$ is isotropic and is fixed by the induced
residue action. Its inverse image $M$ is a positive definite
even unimodular lattice with exactly 72 roots, of root system
$A_2^{12}$ and minimum squared norm two.
The original residue element $r$ is an integral isometry of $M$.
It fixes four of these $A_2$ components pointwise and acts as
a coordinate three-cycle on each of the other eight.

The exact intersection and index relations are
\[
 \begin{gathered}
 M\cap\Lambda_j=J,\qquad [M:J]=[\Lambda_j:J]=9,\\
 M\cap N=L_0:=\{x\in N:\beta(x)=0,\
                                (\varrho,x)\equiv0\pmod2\},\\
 [M:L_0]=[N:L_0]=6.
 \end{gathered}
 \tag{RF33}
\]
If $R_M$ denotes the lattice generated by all its roots, then
\[
 R_M\cong A_2^{\perp12},\qquad
 \det R_M=3^{12}=531441,\qquad[M:R_M]=3^6=729.
 \tag{RF34}
\]
\end{theorem}

\begin{proof}
Substitution of (RF31) in (RF24) gives the numerator
$2a(-b)+b(-a+b)+b(-b)=-3ab=0$ in $\mathbb F_3$.
Thus $P_\infty$ is isotropic.
The matrix (RF27) maps its vector with parameters $(a,b)$ to
\[
 (a+b,b,-a,-b),
 \tag{RF35}
\]
which is again in $P_\infty$, now with parameters $(a+b,b)$.
The images of $p,q$ in (RF23) are
$(1,0,-1,0)$ and $(0,1,1,-1)$, a basis of $P_\infty$.
Consequently (RF32) is its exact inverse image under
$J^\vee\to J^\vee/J$, and $[M:J]=9$.
Isotropy gives even norms on $M$; polarization gives integral
pairings. Its determinant is $\det J/9^2=1$, and its
positive definiteness and rank are inherited from the unchanged
ambient real space. Equation (RF35), together with $rJ=J$,
proves $rM=M$.
Each of $P_0,P_1,P_2$ intersects $P_\infty$ in zero, as follows
by equating their coordinates in (RF26) and (RF31).
Taking inverse images proves $M\cap\Lambda_j=J$.

For its relation with $N$, (RF19) gives the exact expression
\[
 q=ru_*-q_1-\kappa_1 .
 \tag{RF36}
\]
Therefore
\[
 N+M=N+\mathbb Z(\varrho/2)+\mathbb Zq_1.
 \tag{RF37}
\]
The two indicated cosets modulo $N$ have orders two and three.
For the first, pair with $\alpha$; for the second, use (RF10).
Subgroups of coprime orders intersect trivially, so the quotient
in (RF37) has order six.
Also $p\in N$ has $\beta(p)=0$ and
$(\varrho,p)=12-(-4)=16$, hence image $(4,0)$ in (RF15).
Thus $p+I$ has order three.
The vector $3q$ lies in $J$ because $q+J$ has order three,
and lies in $N$ by (RF36) and $3q_1\in N$.
It is therefore in $N\cap J=I$.
We also have $6v\in I$.
Write an element of $M$ using generators $I,3v,p,q$ and reduce
the coefficient of $3v$ modulo two and of $q$ modulo three.
If this element is in $N$, their two coprime-order cosets in
(RF37) force both reduced coefficients to vanish.
Hence
\[
 M\cap N=I+\mathbb Zp.
 \tag{RF38}
\]
Its image in (RF15) is precisely $2(\mathbb Z/6)\oplus0$,
which proves the displayed formula for $L_0$ and
$[N:L_0]=6$. Since $\det N=\det M=1$, the equal covolumes
give $[M:L_0]=6$, proving (RF33).

We next prove that the following root count is exhaustive.
Every vector in $N+M$ has one of the six forms
\[
 x+a\varrho/2-bq_1,\qquad
                   x\in N,\ a\in\{0,1\},\ b\in\{0,1,2\}.
 \tag{RF39}
\]
When $a=0$, the vector is in $R^\vee$.
Every norm-two vector of this particular $R^\vee$ is already
a root of $R$.
Indeed, the $A$ class minima are
$0,5/6,4/3,3/2,4/3,5/6$, and the three nonzero $D$ classes
have minimum one. In a nonzero class of $R^\vee/R$ with
minimum at most two, the possible minima are exactly among
\[
 \{5/6,1,4/3,3/2,5/3,11/6\}.
 \tag{RF40}
\]
This exhausts them because at most two $A$ classes can be
nonzero under the bound; if two are nonzero, each has minimum
$5/6$; with a nonzero $D$ class there is at most one nonzero
$A$ class and its minimum must be $5/6$.
No number in (RF40) is congruent to two modulo $2\mathbb Z$.
All norms within a fixed discriminant class differ by an even
integer, so no such class has a norm-two vector.
This proves the claim about $R^\vee$.

All original roots have $\beta=0$, so (RF33) says that those
lying in $M$ are exactly the ones with even $\varrho$ height.
In each original $A$ block they are
$e_i-e_j$ with $i\ne j$ and $i\equiv j\pmod2$.
These give two orthogonal $A_2$ systems per block, on the
coordinate triples $\{0,2,4\}$ and $\{1,3,5\}$, and 48 roots
in total. In $D$ they are the eight roots
$\pm e_1\pm e_3,\ \pm e_2\pm e_4$.
They form four pairwise orthogonal $A_1$ systems.
This accounts for exactly 56 roots with $a=0$ in (RF39).

For $a=1$, absorb the $A$ dual vector $-bq_1$ into the
marked $A$ discriminant classes.
If the original class of $x$ is $(t_1,t_2,t_3,t_4;y)\in H$,
the shifted class in each $A$ is $t_i-4b$.
The exact half-Weyl-shift minima proved in the neighbor
calculation are
\[
 \min_{w\in t\lambda+A}|w+\varrho_A/2|^2
 =\begin{cases}3/8& t\equiv0\pmod3,\\17/24&t\not\equiv0\pmod3,\end{cases}
 \qquad
 \min_{w+D=y}|w+\varrho_D/2|^2=1/2.
 \tag{RF41}
\]
For completeness the first line follows by choosing centered
quarter coordinates when $t=0,3$; for other $t$, the nearest
twelfth coordinates have square sum $13/24$ and coordinate
sum $\pm1$, and restoring zero sum costs $1/6$.
The second line follows from the integer coordinate distances
$(1/2,0,1/2,0)$ and half-integer distances
$(0,1/2,0,1/2)$, with their ties meeting every $D$ marking.
Thus the total minimum in (RF39) with $a=1$ is at least two.
Equality requires $t_i-4b\equiv0\pmod3$ in every component.
Since $t_i\equiv2z_i\pmod3$, this says
$z_i=2b$ for all four coordinates.
For $b=1,2$ this is excluded by the exact code (RF8).
For $b=0$ it requires $z=0$.
The eight binary codewords and the two $D$ ties then give
exactly sixteen equality vectors $y=x+\varrho/2$.
Every one belongs to $M$: the norm equation gives
\[
 2=x^2+(\varrho,x)+21,\qquad
                         (\varrho,x)=-19-x^2,
 \tag{RF42}
\]
which is odd because $N$ is even.
Also $\beta(x)=0$ because $z=0$.
Therefore $x+3\alpha\in L_0$ and
$y=(x+3\alpha)+3v\in M$.
This proves exhaustively that $M$ has exactly $56+16=72$ roots.

Here are the full root coordinates and the resulting exact
orthogonal decomposition.
Let $\zeta_i$ be supported in the $i$th original $A$ block
with entries
\[
 \zeta=(1,-1,1,-1,1,-1)/4,\qquad \zeta^2=3/8.
\]
In the following table every $Z_y$ is an actual sum of those
four mutually orthogonal vectors, and every $\delta_y$ is
an original $D$ root:
\[
 \begin{array}{c|l|l}
 y&Z_y&\delta_y\\\hline
 00&\zeta_1+\zeta_2+\zeta_3+\zeta_4&e_1+e_3\\
 10&-\zeta_1-\zeta_2+\zeta_3+\zeta_4&e_1-e_3\\
 01&-\zeta_1+\zeta_2-\zeta_3+\zeta_4&e_2+e_4\\
 11&\zeta_1-\zeta_2-\zeta_3+\zeta_4&e_2-e_4
 \end{array}.
 \tag{RF43}
\]
The sixteen new roots are exactly
\[
 \{\epsilon Z_y+\eta\delta_y/2:
                   y\in\mathbb F_2^2,\ \epsilon,\eta\in\{1,-1\}\}.
 \tag{RF44}
\]
To verify the signs from the original glue, the $A$ half-shift
minimizer for binary coordinate $u_i=0$ is $\zeta_i$ and
for $u_i=1$ it is $-\zeta_i$.
For $y=00$ the two codewords are $0000,1111$; for $10$ they
are $1100,0011$; for $01$ they are $1010,0101$; and for $11$
they are $0110,1001$.
These give respectively the four opposite pairs $\pm Z_y$
in (RF43).
The two shifted $D$ minimizers in these exact original
markings are respectively
$\pm(e_1+e_3)/2,\pm(e_1-e_3)/2,
 \pm(e_2+e_4)/2,\pm(e_2-e_4)/2$.
This proves (RF44) with every sign and marking retained.

The four $Z_y$ are mutually orthogonal and have squared norm
$3/2$, as follows by dotting their four displayed sign rows.
The $\delta_y$ are mutually orthogonal of squared norm two.
The $Z_y$ are orthogonal to the entire original $D$ space,
and to the eight old $A_2$ systems because their coordinates
are constant on each parity triple.
For each $y$, the six roots
\[
 \{\pm\delta_y\}
       \ \cup\ \{\epsilon Z_y+\eta\delta_y/2:\epsilon,\eta=\pm1\}
 \tag{RF45}
\]
are an $A_2$ system: take simple roots
$\delta_y$ and $Z_y-\delta_y/2$, whose Gram matrix is
$\left(\begin{smallmatrix}2&-1\\-1&2\end{smallmatrix}\right)$;
their sum is $Z_y+\delta_y/2$.
The four such systems are mutually orthogonal and orthogonal
to the eight old systems. They exhaust the 72 proved roots.
Thus their root lattice is $A_2^{\perp12}$ of determinant
$3^{12}$, and unimodularity of $M$ gives the index $3^6$.
This proves (RF34) and the Coxeter number three of this
rooted Niemeier lattice directly from its twelve components.

Finally $r$ fixes each $\zeta_i$ because it permutes coordinates
within their parity triples; it fixes each $\delta_y$ because
it is the identity on $D$. Hence it fixes the four systems
(RF45) pointwise. On each of the eight old $A_2$ systems it
cyclically permutes its three coordinate positions, giving its
exact order-three Coxeter action.
This completes the asserted classification of the fixed
fourth extension and its precise relation to the three Leech
lattices through $J$.
\end{proof}
